\chapter{Discussion}

Throughout this thesis, we aimed to understand whether and how AI safety evaluations can measure safety discrepancies based on differing user demographics. This was motivated by the suspicion that given known safety failures under, for example, low-resource languages~\cite{yong2023low-resource-ec3, deng2023multilingual-1fc} and evident demographic biases in the models \cite{guo2024biaslargelanguagemodels, hofmann2024ai-e4f}, safeguards may not protect all users equally.

But in order to assess this, as we termed it, \textit{safety discrimination}, one requires evaluations able to explicitly uncover it in the first place. In Chapter \ref{chap:evaleval} we found that safety evaluations centred around the user as victim or adversary in almost no case deliberately evaluated risks against diverse user profiles.

In response to this, we went on to investigate best practices for multilingual evaluations through the case study of sycophancy in Chapter \ref{chap:translations}. Through our own experiment, we found that faithful translation of a diversity of languages and dialects is challenging but not impossible. We also discussed how prior work has investigated in depth different translation approaches along the Pareto frontier between cost and quality \cite{ahuja2023mega-415, xuan2025mmlu-prox-307, peter2025mind-728}.

Moving beyond demographic clues through language, in Chapter \ref{chap:contextualising}, we explored what it takes to bring rich user context into an evaluation. First, we integrated user context in the LLM-judge itself and found that through this context-aware evaluation approach we could uncover safety risks for medium- and high-vulnerability users that a context-blind judge would underestimate. Next, we added user context through the user prompt itself, simulating self-disclosure of relevant demographic information. Using the same evaluation judges as in the previous case study, we found that self-disclosure of user demographics narrows the gap between context-blind and context-aware safety scores for high-vulnerability users but never closes it. This means that context-aware judging should form a core part of future evaluations covering demographic diversity. Lastly, we experimented with memory-based contextualisation as a realistic context source in frontier models. While we could not find significant differences in safety scores stratified by gender or income level, we emphasise that this absence of evidence is much more likely due to our study being very simplistic and heavily underpowered in data points than grounds to draw conclusions on the outcomes.

In the final experimental chapter, Chapter \ref{chap:probes}, we observed that linear probes built to detect social sycophancy \cite{cheng2025social-928} have differing transferability depending on the training data language. Probes trained on high-resource languages such as English, German or Russian only slightly degrade in performance when tested to monitor sycophancy in mid-resource languages. However, the inverse does not hold true, and a probe trained on mixed-language data achieves similarly low accuracy. This underscores that even safeguarding techniques such as monitors might suffer from a lack of demographic generalisation.

Coming back to where we started – the question about methodological and behavioural safety discrimination – we now want to discuss how our results answer and serve each of them, where the limitations of our investigations are and what we suggest future work should cover.

\section{Methodological Safety Discrimination}
With methodological safety discrimination, we described the case when the methods to measure and provide safety are themselves not designed with a pluralistic approach \cite{sorensen2024roadmap} to demographics in mind.

In Chapter \ref{chap:evaleval}, we established that this indeed seems to be the case for user safety evaluations. Out of the 23 evaluations we looked at, ten did not even mention missing demographic diversity as a limitation. We recognise that our selection of evaluations provides only limited breadth and quantity, which we traded off against quality through peer review, and thus there might be evaluations for very specific demographics that our analysis missed. In such cases, the question remains to what degree frontier AI developers are using those non-standard evaluations in their model evaluations.
Specifically for languages, established protocols for multilingual approaches exist, but can quickly get expensive \cite{ahuja2023mega-415, xuan2025mmlu-prox-307}.

For inclusion of other demographics in evaluations, approaches are sparse, which is why we experimented ourselves with novel methods in Chapter \ref{chap:contextualising}. While those experiments establish preliminary frameworks, they are limited by uncertainties around the validity of the LLM-judge \cite{gu2025surveyllmasajudge, schroeder2025trustllmjudgmentsreliability}, the model selection and limited sample sizes. We encourage future work to iterate and improve on our frameworks to alleviate those limitations.

But even with methods available, a larger question overarches this discussion: the sample space of possible combinations of demographics is vast and evaluations covering every possible case are hence impossible. Nevertheless, we advocate for at least the following approaches to deal with this problem: first, as we have done ourselves in Chapter \ref{chap:contextualising}, evaluations should focus on prioritising demographics by vulnerability. While there will always be exceptions, assuming high-vulnerability users as the default will, if at all, overestimate risks but not harm users with lower vulnerability. While evaluations stratified by demographics or vulnerability may provide the most accurate harm estimates, we agree with
prior work critiquing the construct invalidity of many safety evals \cite{raji2021everything, jacobs2021measurement, weidinger2025evaluationscience},
that at a bare minimum, evaluations should clearly define the scope of what the evaluation is aiming to measure and for which user population to prevent generic safety claims that do not hold up for specific user groups.

On the front of safety discrimination in the methods of safeguards themselves, we only have limited evidence from our study of generalisation of linear probes monitoring for social sycophancy across languages. Our result of limited cross-lingual generalisation is based only on our investigation of a single 14B model and social sycophancy as monitoring target specifically. Whether linear probes for monitoring show concerning performance differences dependent on language or other demographics across model families and sizes, safety-relevant behaviours, and a wider variation of demographics, we leave to future work. In our limited result however, we found that probes trained on English data transfer well across languages, which is a good sign under the assumption that probes used in deployment \cite{mckenzie2025detecting-612} are primarily trained on English data.

However, this finding again may not hold true in general and therefore dedicated generalisation evaluations should be conducted rather than relying on unverified assumptions. Future work should also stress-test other safeguarding techniques on their cross-demographic generalisation.

Despite the limitations of our work in uncovering the \textit{extent} of methodological safety discrimination, we found evidence for the \textit{presence} thereof, inviting future work to further study where methods fall short and how to mitigate those failures.


\section{Behavioural Safety Discrimination}
Given the limited number of methods that would allow us to uncover behavioural safety discrimination as we established in the previous section, the question whether models themselves behave in a discriminatory way in the safety of their responses towards certain user groups stays open. While for translations, methods exist and yet the actual evaluations are largely missing, for other demographic factors apart from language, we may need novel methods in the first place. We aimed to provide three exploratory approaches through our case studies in Chapter \ref{chap:contextualising}. While we could not validate our LLM-judge against human expert annotations, the approach of a context-aware LLM-judge uncovered otherwise hidden risks and followed consistent directional patterns throughout the experiments, allowing for a careful recommendation to follow up on this approach. On context disclosure at the user front, where we explored prompt- and memory-based approaches, we recommend that future work should orient and align with what context disclosure looks like at the frontier \cite{eloundou2025firstperson}.

In particular, we are excited to see more evaluations leveraging the user memory, which is itself easily exportable from chat interfaces like ChatGPT or Claude. Realistic context-dependent evaluations may hence be possible through memory donations paired with collecting key demographic information about the user to investigate memory-based contextualisation together with context-aware judging, which we have not done in this work.

In general, we want to note that as the frontier of AI progresses, models may generalise better anyway, or external evaluations may only be partially feasible or only partially provoke change inside the labs, especially given the exact design of safeguards largely stays hidden from the public.

Besides evaluations, public model specs~\cite{openai2025modelspec, anthropic2026constitution} may be a feasible target to anchor demographic inclusiveness and may be usable to hold companies accountable to their commitment.

Finally, there remains the question to what extent the appropriateness of behaviour is culturally dependent \cite{kirk2024prism, aroyo2023dices}. In our case study of social sycophancy, this may mean that in some cultures expressed through the respective language, a response classified as sycophantic in one culture may be appropriately polite in another. For this, cultural calibration studies may be an important step towards building datasets that ground demographics in cultural context.



\section{Limitations}
Besides the limitations named for each individual case study, it is important to also look at the limitations of this thesis as a whole.
While discussing a broader phenomenon, we underscored the issues through case studies, each focussed on narrow risks and a small sample of stratified demographic variables. Furthermore, all our experiments cover only a limited number of models, mostly small ones which were cheaper to run. This leaves us with scattered findings and points of evidence, but insufficient ground to finally answer whether safety discrimination is a broadly occurring problem. Nevertheless, we hope that this thesis can point future work to drive progress on this question by developing vulnerability-prioritised evaluations and conducting empirical studies of generalisation of safeguards close to the frontier.

